\documentclass[11pt]{article}
\usepackage[margin=1in]{geometry}
\usepackage{parskip}
\usepackage[colorlinks=true,allcolors=blue]{hyperref}

\begin{document}

\pagestyle{empty}

\noindent Hikaru Wakaura and Taiki Tanimae\\
QIRI (Quantum Integrated Research Institute Inc.)\\
Tokyo 107-0061, Japan\\
\href{mailto:h.wakaura@qiri.co.jp}{h.wakaura@qiri.co.jp} \quad ORCID: 0000-0001-8381-8323

\vspace{1em}
\noindent July 10, 2026

\vspace{1em}
\noindent To the Editors, \emph{PRX Quantum}

\vspace{1em}

\noindent \textbf{Re: Submission of ``Scrambling, not athermality, protects
quantum information under generic measurement: chaotic states outperform
quantum many-body scar codes''}

\vspace{0.5em}

On July~10, 2026, we submit for your consideration the manuscript above (research
article, two authors) together with Supplemental Material (one PDF: eight
sections, three figures, one table) and an openly available simulation package
reproducing every result.

\vspace{0.5em}

% Paragraph 1: the surprise (overturns a common intuition) + what it reveals
A widely shared intuition holds that quantum many-body scars---non-thermal,
weakly entangled, unusually long-lived---should make especially good quantum
memories. We show that this intuition is wrong, and in doing so identify what
actually makes a many-body code robust to measurement. Computing the
reference-qubit coherent information of scar versus thermal codes in a monitored
circuit, we find that generic chaotic states protect information \emph{better}
under every realistic measurement and at every accessible size in the
paradigmatic PXP model---down to the largest code an entire scar subspace can
hold---while an exactly solvable second model confirms the mechanism and shows
that scars gain an advantage only under a single fine-tuned probe. The result
is counterintuitive by construction: the very structure that singles scars out
is precisely what a measurement exploits to read them out.

% Paragraph 2: broad significance (design principle + methodological lesson) — preempts "confirmatory"
Two features give the work reach well beyond scars. First, it yields a general
design principle for measurement-robust quantum memory: protection is set by a
code's spread in the measurement basis, not by any thermodynamic property, and
structured subspaces whose symmetry generators couple the encoded states are
actively harmful. We make this falsifiable---an exactly solvable model isolates a
single control parameter, a Casimir variance, and predicts correctly the unique
fine-tuned measurement under which scars can be rescued. Second, of immediate
practical value to the many groups now benchmarking codes and channels, we
document a reproducibility trap: comparing against a single natural-looking
reference eigenstate inverted our own initial conclusion, and only ensemble
averaging exposed it as an outlier---a cautionary result we consider broadly
useful in its own right. Together these carry the work past any mere confirmation
of the thermalization--error-correction correspondence: they turn that
correspondence into an operational, testable statement about code \emph{design},
overturn the opposing folklore with a decisive exact-model test, and flag a
pitfall in how such statements are established.

\noindent In brief, relative to prior work: \emph{(i)} the first head-to-head
coding comparison of scar versus thermal subspaces under monitoring, at the
level of channel capacity rather than entanglement or fidelity; \emph{(ii)} a
two-channel protection criterion with an exact single-shot predictor, validated
across 65 code--measurement pairs in two models; \emph{(iii)} a documented
sign-reversing baseline artifact of direct relevance to anyone benchmarking
codes against individual eigenstates.

\vspace{0.5em}

% Paragraph 3: why this journal (broad quantum-info interest, conceptual, reproducible)
We believe \emph{PRX Quantum} is the right venue. The result speaks to the whole
quantum-information community working on monitored dynamics, error correction,
and quantum memory; it is conceptual and general rather than model-specific; and
it is delivered with full reproducibility---an open package regenerates every
figure with one command, and all quantitative claims have been independently
checked against primary sources. We hope the Editors will agree that overturning
a common intuition, extracting a design principle, and surfacing a shared
methodological hazard together meet the journal's bar for significance and
breadth.

\vspace{0.5em}

\noindent We suggest the following referees, all leading experts in
measurement-induced dynamics, quantum error correction, or quantum many-body
scars (institutional contact addresses are publicly listed):

\begin{itemize}\itemsep2pt
  \item Prof.\ Matthew P.~A.\ Fisher, University of California, Santa Barbara ---
        measurement-induced transitions
  \item Prof.\ Ehud Altman, University of California, Berkeley ---
        error correction and monitored circuits
  \item Dr.\ Michael J.~Gullans, NIST / University of Maryland ---
        purification transitions
  \item Prof.\ Vedika Khemani, Stanford University --- random and monitored
        circuits
  \item Prof.\ Adam Nahum, \'Ecole Normale Sup\'erieure, Paris --- entanglement
        dynamics
  \item Prof.\ Zlatko Papi\'c, University of Leeds --- quantum many-body scars
  \item Prof.\ Maksym Serbyn, IST Austria --- quantum many-body scars
  \item Dr.\ Sanjay Moudgalya, Technical University of Munich --- exact scars and
        spectrum-generating algebras
\end{itemize}

\noindent We do not request the exclusion of any referee.

\vspace{1em}

\noindent On behalf of both authors,\\
Hikaru Wakaura\\
QIRI (Quantum Integrated Research Institute Inc.), Tokyo 107-0061, Japan\\
\href{mailto:h.wakaura@qiri.co.jp}{h.wakaura@qiri.co.jp}

\end{document}
